% Section: Recursive Development Systems
\section{Recursive Development Systems}
\label{sec:recursive-development}

% TODO: Expand this section.
%
% Key topics:
% - Definition of recursive development in the context of AI-augmented engineering
% - Distinction between iterative and recursive development patterns
% - How autonomous agents create self-referential feedback loops
% - Benefits and risks of recursive automation
% - The role of scaffolding, code generation, and agentic pipelines
% - Examples of recursive development patterns in modern AI tooling

Recursive development systems are software engineering workflows in which AI agents
participate in their own construction, validation, and refinement cycles. Unlike
traditional iterative development, where human developers drive each cycle, recursive
systems allow autonomous agents to invoke sub-agents, evaluate their own outputs, and
initiate subsequent development passes without direct human intervention at each step.

This recursive structure offers significant productivity advantages: agents can parallelize
work, maintain context across long-horizon tasks, and apply consistent coding standards
without fatigue. However, recursive autonomy introduces compounding risk surfaces.
Each recursive invocation may inherit and amplify errors, assumptions, or misalignments
from prior passes—a phenomenon we term \textit{recursive drift}.

\subsection{Characteristics of Recursive Development}

\begin{itemize}
  \item \textbf{Self-referential invocation}: Agents invoke further agents as part of
    their execution strategy.
  \item \textbf{Compounding context}: Each recursive pass builds on the artifacts of
    prior passes, creating layered dependency chains.
  \item \textbf{Autonomous pacing}: Without external gates, recursive systems can
    progress indefinitely without human review.
  \item \textbf{Emergent complexity}: System behavior becomes increasingly difficult
    to predict as recursion depth grows.
\end{itemize}

\subsection{The Recursion Boundary Problem}

A central challenge in recursive development is determining when and how to interrupt
autonomous execution for human review. Systems that interrupt too frequently lose the
efficiency benefits of automation; systems that interrupt too rarely risk undetected
drift and misalignment.

Reflector addresses this challenge through \textit{milestone synchronization}, described
in Section~\ref{sec:milestone-synchronization}, which defines explicit recursion
boundaries tied to semantic milestones rather than arbitrary time or iteration counts.
